% EM structure and teaching revision 20260928
\subsection{分层介质：同电压还是同自由电荷}
对线性、各向同性、无明显边缘效应的平板结构，介质分层方向决定等效连接。若两层沿电场方向依次排列，厚度为 $d_1,d_2$，界面无自由电荷，则法向位移 $D=Q/S$ 在两层中相同，电压为
\[
 U=E_1d_1+E_2d_2=\frac QS\left(\frac{d_1}{\varepsilon_1}+\frac{d_2}{\varepsilon_2}\right),\qquad
 C=\frac S{d_1/\varepsilon_1+d_2/\varepsilon_2}.
\]
若介质左右并排、共同跨接两块极板，则两部分电压相同，电荷相加，得到
$C=(\varepsilon_1S_1+\varepsilon_2S_2)/d$。不能仅凭“各占一半”选择串联或并联。
取 $d_1=d_2=d/2$、$\varepsilon_1=\varepsilon_0$、$\varepsilon_2=4\varepsilon_0$，沿场分层得到 $C=1.6\varepsilon_0S/d$，等面积并排则为 $2.5\varepsilon_0S/d$。低介电常数层中的场更强，也更可能先达到击穿阈值。
上述计算假定介质响应线性且几何保持不变。若要求介质受力，还要说明电源是否连接，并把电源的能量交换纳入功的计算。

